% ============================================================
% Table: Same-day Calibration on SPY (2022-09-02, 100 contracts)
% ============================================================
\begin{table}[htbp]
  \centering
  \caption{Same-Day Calibration Performance on SPY Options (2022-09-02, $M=100$ contracts)}
  \label{tab:spy_sameday}
  \setlength{\tabcolsep}{7pt}
  \begin{tabular}{lccc}
    \toprule
    \textbf{Metric} & \textbf{Proposed Method} & \textbf{Zhang (2025)} & \textbf{Better} \\
    \midrule
    \multicolumn{4}{l}{\textit{Native model output}} \\
    \quad IV MRE & \textbf{1.56\%} & 5.01\% & Proposed \\
    \quad Price MRE & n/a & \textbf{13.36\%} & --- \\
    \midrule
    \multicolumn{4}{l}{\textit{QuantLib re-pricing with calibrated parameters}} \\
    \quad IV MRE & \textbf{8.15\%} & 9.18\% & Proposed \\
    \quad Price MRE & 18.10\% & \textbf{16.93\%} & Zhang \\
    \midrule
    \multicolumn{4}{l}{\textit{Calibration setup}} \\
    \quad Objective function & Vega-weighted IV MSE + Feller penalty & Normalised price MSE & --- \\
    \quad Feller constraint enforced & Yes & No & --- \\
    \quad Number of restarts & 10 & 10 & --- \\
    \quad Max steps per restart & 300 & 300 & --- \\
    \bottomrule
  \end{tabular}
  \begin{tablenotes}
    \small
    \item \textit{Note:} IV MRE $= \frac{1}{M}\sum_{i=1}^{M} |\hat{\sigma}^{\text{IV}}_i - \sigma^{\text{IV},\text{mkt}}_i|\,/\,|\sigma^{\text{IV},\text{mkt}}_i|$. The proposed method calibrates directly in implied-volatility space, so it does not natively output a price-space MRE. For Zhang~(2025), the native IV rows are obtained by converting DDN-predicted prices to Black--Scholes implied volatilities on the 100-contract SPY chain used in the price-space notebook. QuantLib rows re-price the same chain using the calibrated Heston parameters.
  \end{tablenotes}
\end{table}


% ============================================================
% Table: Cross-Day Calibration on SPY (9/1 params → 9/2 chain)
% ============================================================
\begin{table}[htbp]
  \centering
  \caption{Cross-Day Calibration Performance on SPY Options (Parameters from 2022-09-01 Applied to 2022-09-02, $M=100$ contracts)}
  \label{tab:spy_crossday}
  \setlength{\tabcolsep}{7pt}
  \begin{tabular}{lccc}
    \toprule
    \textbf{Metric} & \textbf{Proposed Method} & \textbf{Zhang (2025)} & \textbf{Better} \\
    \midrule
    \multicolumn{4}{l}{\textit{Native model output on the 2022-09-02 chain}} \\
    \quad IV MRE & \textbf{1.98\%} & 3.71\% & Proposed \\
    \quad Price MRE & n/a & \textbf{10.97\%} & --- \\
    \midrule
    \multicolumn{4}{l}{\textit{QuantLib re-pricing with cross-day parameters}} \\
    \quad IV MRE & \textbf{9.23\%} & 10.79\% & Proposed \\
    \quad Price MRE & \textbf{18.89\%} & 20.70\% & Proposed \\
    \midrule
    \multicolumn{4}{l}{\textit{Reference: same-day native IV fit on 2022-09-02}} \\
    \quad IV MRE (same-day) & 1.56\% & 5.01\% & --- \\
    \midrule
    \multicolumn{4}{l}{\textit{Change in native IV MRE relative to same-day}} \\
    \quad $\Delta$ IV MRE & $+$0.42\,pp & $-$1.30\,pp & See note \\
    \bottomrule
  \end{tabular}
  \begin{tablenotes}
    \small
    \item \textit{Note:} ``pp'' denotes percentage-point change in MRE. The proposed method is calibrated on 2022-09-01 and then evaluated on the 2022-09-02 option chain without re-calibration. For the Zhang baseline, the native IV metric is derived from price outputs after Black--Scholes inversion, so its same-day and cross-day IV MRE need not move monotonically with the price-space training objective. The more direct cross-day comparison is that the proposed method yields lower next-day IV error and lower QuantLib re-pricing error on the evaluation day.
  \end{tablenotes}
\end{table}
